\subsection{W 的 Coxeter 矩阵与 Coxeter 图}

设 C 是一个室，\(S = S(C)\) 为关于 C 的墙的正交反射所成的集合，\(M = (m(s, s'))\) 为 Coxeter 系统 (W, S) 的 Coxeter 矩阵（《代数》第 IV 章 §1 第 9 号）；回顾 \(m(s, s')\) 是 W 中元素 \(ss'\) 的阶（有限或无限）（其中 \(s, s' \in S\)）。若 \(C'\) 是另一个室，则满足 \(w \in W\) 且 \(w(C) = C'\) 的唯一元素 w 定义一个双射

\[
s \mapsto f (s) = w s w ^ {- 1}
\]

从 S 到 \(S' = S(C')\)，并且有 \(m(f(s), f(s')) = m(s, s')\)。由此，如果 W 通过 \((C, s)\) 作用于由二元组 \(s \in S(C)\) 构成的集合 X，其中 C 是一个室且 \(w.(C, s) = (w(C), wsw^{-1})\)，则 W 在 X 中的每个轨道都与集合 \(\{C\} \times S(C)\) 恰好相交于一点，我们记该点为 \((C, s_i(C))\)。因此，若 I 是轨道集且 \(i, j \in I\)，数值 \(m_{ij} = m(s_i(C), s_j(C))\) 与室 C 的选择无关。矩阵 \(M(W) = (m_{ij})_{i,j \in I}\) 是一个 Coxeter 矩阵，称为 W 的 Coxeter 矩阵。与 M(W) 相关的 Coxeter 图（《代数》第 IV 章 §1 第 9 号）称为 W 的 Coxeter 图。

设 C 是一个室。对于任意 \(i \in I\)，记 \(\mathrm{H}_{i}(\mathrm{C})\) 为 C 的满足下述条件的墙：\(s_{i}(\mathrm{C})\) 是关于 \(\mathrm{H}_{i}(\mathrm{C})\) 的反射；并记 \(e_{i}(\mathrm{C})\) 为与 \(\mathrm{H}_{i}(\mathrm{C})\) 正交、且位于 \(\mathrm{H}_{i}(\mathrm{C})\) 与 C 同一侧的单位向量。映射 \(i \mapsto \mathrm{H}_{i}(\mathrm{C})\) 称为 C 的规范编号墙族。

命题 3. 设 \(\mathbf{C}\) 是一个室，且 \(i, j \in \mathbf{I}\) 满足 \(i \neq j\)。置 \(s_i = s_i(\mathbf{C})\)、\(\mathbf{H}_i = \mathbf{H}_i(\mathbf{C})\)、\(e_i = e_i(\mathbf{C})\)，并类似地定义 \(s_j\)、\(\mathbf{H}_j\) 和 \(e_j\)。

\begin{enumerate}
\def\labelenumi{(\roman{enumi})}
\item
若 \(\mathrm{H}_i\) 与 \(\mathrm{H}_j\) 平行，则 \(m_{ij} = \infty\) 且 \(e_i = -e_j\)。
\item
若 \(\mathrm{H}_i\) 与 \(\mathrm{H}_j\) 不平行，则 \(m_{ij}\) 有限，且
\end{enumerate}

\[
\left(e _ {i} \mid e _ {j}\right) = - \cos (\pi / m _ {i j}).\tag{4}
\]

\begin{enumerate}
\def\labelenumi{(\roman{enumi})}
\setcounter{enumi}{2}
\tightlist
\item
有 \((e_i|e_j)\leqslant 0\)。
\end{enumerate}

若 \(\mathrm{H}_i\) 与 \(\mathrm{H}_j\) 平行，则 \(s_i s_j\) 是一个平移（§2，第 4 号，命题 5），所以 \(m_{ij} = \infty\)。此外，或者 \(e_i = e_j\)，或者 \(e_i = -e_j\)。现在，闭包 C 中存在一点 \(a\)（分别为 \(a'\)），它属于 \(\mathrm{H}_i\)（分别属于 \(\mathrm{H}_j\)），但不属于 \(\mathrm{H}_j\) \((\text{resp. } \mathrm{H}_i)\)。于是 \((a' - a|e_i) > 0\) 且 \((a - a'|e_j) > 0\)，这排除了 \(e_i = e_j\) 的情形，并证明了 (i)。

现在假设 \(H_{i}\) 与 \(H_{j}\) 不平行。取 \(a \in H_{i} \cap H_{j}\) 为原点，并通过双射 \(t \mapsto a + t\) 将 T 与 E 识别。令 V 为与 \(H_{i} \cap H_{j}\) 正交且经过 a 的平面。置 \(\Gamma = V \cap D_{H_{i}}(C) \cap D_{H_{j}}(C)\)（其中 \(D_{H}(C)\) 表示由 H 界定且包含 C 的开半空间（§1，第 1 号）），并令 D（分别为 \(D'\)）为 V 中包含于 \(H_{i} \cap V\)（分别为 \(H_{j} \cap V\)）且包含于 \(\Gamma\) 闭包中的半直线。对 V 取适当的定向后，集合 \(\Gamma\) 是 V 中满足下式的开半直线 \(\Delta\) 的并：

\[
0 <   (\widehat {\mathrm{D} , \Delta}) <   (\widehat {\mathrm{D} , \mathrm{D} ^ {\prime}}).
\]

令 \(W'\) 为由 \(W\) 和 \(s_i\) 生成的 \(s_j\) 的子群。对于任意 \(w \in W'\)，超平面 \(w(H_i)\) 和 \(w(H_j)\) 属于 \(\mathfrak{H}\)，包含 \(H_i \cap H_j\)，且不与 C 相交。因此，它们不与 \(\Gamma\) 相交（§1，第 5 号，命题 10）。于是，§2 第 5 号命题 7 的推论蕴含 (ii)。

最后，断言 (iii) 直接由 (i) 和 (ii) 得出，因为 \(m_{ij} \geqslant 2\) 对 \(i \neq j\) 成立。

注. 公式 (4) 实际上对所有 \(i, j \in I\) 都成立：事实上，当 \(\pi / m_{ij} = 0\) 时，\(m_{ij} = \infty\)，而当 \(i = j\) 时，\(m_{ij} = 1\) 且 \((e_i | e_j) = 1\)。

